%! Tex main = MarquezSalazarBrandon-PropuestaDeTesis.tex

\section{Introducción}

El uso de herramientas computacionales ha tenido un impacto
significativo dentro de la investigación científica \cite{Du_2004}.
Desde el desarrollo de sistemas simples de tareas de
automatización y gestión de datos \cite{WOS:A1985ACW9900019},
hasta modernas aproximaciones de inteligencia artificial basada
en tensores \cite{Yenduri_2024, Wilson_2014}.

\subsection{Sistemas inspirados en la naturaleza}

Dentro de este ámbito se encuentran los \textit{nature-inspired methods} 
(NIM) o métodos inspirados en la naturaleza. Dichos métodos
se basan en el comportamiento de seres vivos, procesos físicos,
interacciones protéicas, entre otros \cite{YANG20203, WOS:000975679700001}.

El objetivo principal es la creación de un conjunto de reglas
simples que puedan describir un comportamiento más complejo;
generalmente estos métodos permiten la exploración dentro del
espacio de búsqueda \cite{YANG20203, Yang2010, Yang2017, Lones_2014}.

Este tipo de implementaciones permiten encontrar soluciones en
modelos cuyo comportamiento es desconocido. Tal como se
menciona en \cite{Eiben_2015}, la optimización consiste en la
búsqueda de valores de entrada óptimos para un modelo.

\subsubsection{Heurísticas, metaheurísticas y NIM}

Una diferenciación necesaria dentro de esta área es el uso de
heurística, metaheurística y método inspirado en la naturaleza.

Mientras las heurísticas son métodos que permiten la búsqueda
de posibles soluciones, su diferencia con las metaheurísiticas
es que estas últimas tiene un comportamiento más complejo que
puede equipararse al aprendizaje \cite{Lones_2014, Yang2017}.

Por su parte, un NIM suele tomar reglas del proceso a imitar y
crear un conjunto de pasos a seguir de forma iterativa a fin de 
producir comportamientos complejos, muchas veces adaptivos e
\say{\textit{inteligentes}} \cite{Holland_1992, Yang2017, Eiben_2015}.

\subsubsection{Exploración y explotación}

Generalmente los modelos en los que se utilizan los NIM definen
espacios no lineales donde métodos deterministas que,
generalmente, requieren de gradientes o de conocimiento previo
del comportamiento del sistema, suelen fallar o quedarse en
óptimos locales \cite{WOS:000975679700001, Yang2017}.

El comportamiento de los modelos más deterministas permiten
afinar un punto dentro de los límites propuestos en el espacio
de búsqueda, precisión que, generalmente, los modelos
metaheurícos pierden a cambio de la perspectiva global \cite{Yang2017}.

Dicha búsqueda o mapeo suele llamarse exploración, mientras que
el refinamiento determinístico de puntos óptimos es llamado
explotación.

\subsection{Teoría de cuerdas}

La teoría de cuerdas en una rama de la física que propone a las
partículas fundamentales del universo  como cuerdas de energía
unidimensionales en vez de puntos adimensionale, lo que ha
permitido obtener modelos matemáticos más estables \cite{Becker2007}.

Además de lo anterior, la teoría de cuerdas apunta a la
posibilidad de definir la gravedad en términos cuánticos \cite{Rickles_2014}.

Entre los principales objetivos de la investigación en teoría
de cuerdas se encuentran la construcción de un modelo que
nos permita describir las partículas de nuestro universo.

\subsubsection{El \textit{landscape} y el \textit{swampland}}

La teoría de cuerdas por sí misma es un marco teórico que
permite modelar distintos universos. Cada universo se rige
por un conjunto de leyes que se derivan de las diferentes
interacciones entre sus cuerdas.

Existen un conjunto de teorías y conjeturas que han permitido
dirigir la investigación, sin embargo, el proceso sigue siendo
exhaustivo matemáticamente y topológicamente \cite{Rickles_2014}.

Debido a lo anterior, la teoría de cuerdas ha enfrentado uno de
sus mayores retos: el \textit{landscape}. Este se refiere a la gran
cantidad de universos posibles.

Por su parte, el \textit{swampland} se refiere al conjunto de universos
verosímiles a un modelo estable del universo que, sin embargo,
son inestables o inconsistentes \cite{Vafa2005}.
